\documentclass[10pt]{article}
\usepackage[margin=0.78in]{geometry}
\usepackage{amsmath,amssymb}
\usepackage[colorlinks=true,linkcolor=blue,urlcolor=blue,citecolor=blue]{hyperref}
\setlength{\parindent}{0pt}
\setlength{\parskip}{0.42em}

\begin{document}
\begin{center}
{\Large\bfseries Literature Status: Talagrand's Problem 1 for Boolean Convolution}\par
\smallskip
{\small fa\_banach\_001 literature-status packet -- 19 July 2026}
\end{center}

\textbf{Status.} \texttt{literature\_already\_answered} for Problem 1 only.
The sharp quantitative Problem 2 remains open in the supporting source.

\section*{Original source and question}

Gozlan, Li, Madiman, Roberto, and Samson~\cite{GLMRS} state Talagrand's two
Boolean heat-semigroup conjectures in the Introduction on PDF page 2.  With
\(\mu\) uniform on \(\{-1,1\}^n\), \(P_\tau\) the heat semigroup, and
\(f\geq0\) normalized by \(\lVert f\rVert_1=1\), Problem 1 asks whether
\[
 \lim_{\eta\to\infty}
 \eta\sup_n\sup_f \mu\{P_\tau f\geq\eta\}=0
 \qquad(\tau>0).
\]
Problem 2 asks for the stronger sharp estimate
\[
 \eta\sup_n\sup_f \mu\{P_\tau f\geq\eta\}
 \leq \frac{c_\tau}{\sqrt{\log\eta}}.
\]
The source says on the same page that both conjectures were open.

\section*{Supporting answer and identification}

Chen's Theorem 1~\cite{Chen}, on PDF page 2, proves that for
\(\tau>0\), \(n\geq1\), \(\eta>e^3\), and nonnegative \(f\),
\[
 \mu\{P_\tau f>\eta\lVert f\rVert_1\}
 \leq
 \frac{c}{(1-e^{-\tau})^2}
 \frac{(\log\log\eta)^{3/2}}{\eta\sqrt{\log\eta}}.
\]
Multiplication by \(\eta\) and passage to the limit immediately gives the
limit in Problem 1.  More importantly for provenance, the paragraph directly
after Theorem 1 explicitly names Talagrand's Problem 1 and states that the
theorem resolves it.  Thus the supporting author knowingly answers the same
problem; the relation is not merely inferred in this packet.

\section*{Scope limitation}

The factor \((\log\log\eta)^{3/2}\) still diverges, so Theorem 1 does not give
the sharp Problem 2 estimate.  The supporting paper itself presents its result
as resolving the conjecture only up to this dimension-free loss.  Accordingly,
this packet records a literature answer to Problem 1 only and makes no claim
that Problem 2 is solved.

\begin{thebibliography}{9}
\bibitem{GLMRS}
N. Gozlan, X.-M. Li, M. Madiman, C. Roberto, and P.-M. Samson,
\emph{Log-Hessian and Deviation Bounds for Markov Semi-Groups, and
Regularization Effect in L1}, arXiv:1907.10896v2, 2022.

\bibitem{Chen}
Y. Chen,
\emph{Talagrand's convolution conjecture up to log log via perturbed reverse
heat}, arXiv:2511.19374, Theorem 1, 2025 (current arXiv version consulted
19 July 2026).
\end{thebibliography}

\end{document}
